\par\Needspace{6\baselineskip}\medskip
\hypertarget{native-091-article071}{}\pdfbookmark[2]{71.}{article-071}
\begin{center}71.\end{center}
Die Aberration der Fixsterne und auch derjenige Teil der Aberration der Planeten und Kometen, der allein auf die Bewegung der Erde zurückgeht, entstehen dadurch, dass sich der \emph{Tubus} mit der ganzen Erde bewegt, während ein Lichtstrahl seine optische Achse durchläuft. Der beobachtete Ort eines Himmelskörpers (auch sein scheinbarer oder von der Aberration betroffener Ort genannt) wird durch die Lage der optischen Achse eines so aufgestellten Fernrohrs bestimmt, dass der vom Himmelskörper ausgehende Lichtstrahl auf seinem Weg beide Enden dieser Achse trifft. Diese Lage unterscheidet sich von der wahren Lage des Lichtstrahls im Raum. Wir unterscheiden die beiden Zeitpunkte $t$ und $t'$, zu denen der Lichtstrahl das vordere Ende (die Mitte des Objektivglases) beziehungsweise das hintere Ende (den Brennpunkt des Objektivglases) erreicht. Die Orte dieser Enden im Raum seien zum früheren Zeitpunkt $a$ und $b$ und zum späteren $a'$ und $b'$. Offenbar gibt die Gerade $ab'$ die wahre Lage des Strahls im Raum an, während dem scheinbaren Ort die Gerade $ab$ oder $a'b'$ entspricht (beide Geraden darf man als parallel ansehen). Ferner \SourceLeaf{092}sieht man leicht, dass der scheinbare Ort von der Länge des Tubus unabhängig ist. Der Unterschied in der Lage der Geraden $b'a$ und $ba$ ist die Aberration, wie sie bei Fixsternen auftritt. Ihre Berechnung setzen wir hier als bekannt voraus und übergehen sie. Bei Wandelsternen ist dieser Unterschied jedoch noch nicht die volle Aberration. Der Planet selbst ändert nämlich seinen Ort, während der von ihm ausgegangene Lichtstrahl zur Erde gelangt. Daher entspricht die Lage dieses Strahls nicht dem wahren geozentrischen Ort zur Zeit der Beobachtung. Wir nehmen an, dass der Lichtstrahl, der zur Zeit $t$ in den Tubus eintritt, zur Zeit $T$ vom Planeten ausgegangen ist. Der Ort des Planeten im Raum werde zur Zeit $T$ mit $P$ und zur Zeit $t$ mit $p$ bezeichnet. Schließlich sei $A$ der Ort des vorderen Endes der Tubusachse zur Zeit $T$. Dann gilt
\par $1^{\mathrm{o}}$ Die Gerade $AP$ gibt den wahren Ort des Planeten zur Zeit $T$ an.
\par $2^{\mathrm{o}}$ Die Gerade $ap$ gibt den wahren Ort zur Zeit $t$ an.
\par $3^{\mathrm{o}}$ Die Gerade $ba$ oder $b'a'$ gibt den scheinbaren Ort zur Zeit $t$ oder $t'$ an (der Unterschied der beiden Zeiten kann als unendlich kleine Größe angesehen werden).
\par $4^{\mathrm{o}}$ Die Gerade $b'a$ gibt denselben scheinbaren Ort an, von der Aberration der Fixsterne befreit.

Nun liegen die Punkte $P$, $a$ und $b'$ auf einer Geraden, und die Strecken $Pa$ und $ab'$ sind den Zeitintervallen $t-T$ und $t'-t$ proportional, da sich das Licht mit gleichförmiger Geschwindigkeit bewegt. Wegen der ungeheuren Lichtgeschwindigkeit ist das Zeitintervall $t'-T$ stets sehr klein. Währenddessen darf die Bewegung der Erde als geradlinig und gleichförmig angenommen werden. Somit liegen auch $A$, $a$ und $a'$ auf einer Geraden, und die Strecken $Aa$ und $aa'$ sind ebenfalls den Intervallen $t-T$ und $t'-t$ proportional. Daraus folgt leicht, dass die Geraden $AP$ und $b'a'$ parallel sind und demnach der erste Ort mit dem dritten übereinstimmt.

Die Zeit $t-T$ ergibt sich als Produkt des Abstands $Pa$ mit $493^{\mathrm{s}}$, der Zeit, in der das Licht den mittleren Abstand der Erde von der Sonne durchläuft, den wir als Einheit gewählt haben. In dieser Rechnung dürfen wir für den Abstand $Pa$ auch $PA$ oder $pa$ einsetzen, da der Unterschied ohne Belang sein kann.

Aus diesen Grundsätzen ergeben sich drei Verfahren, um den scheinbaren Ort eines Planeten oder Kometen für eine beliebige Zeit $t$ zu bestimmen. Je nach den Umständen ist das eine oder das andere vorzuziehen.

I. Man ziehe von der gegebenen Zeit die Zeit ab, in der das Licht vom Planeten zur Erde gelangt. So erhält man die reduzierte Zeit $T$, für die der auf die übliche Weise berechnete wahre Ort mit dem scheinbaren Ort für $t$ übereinstimmt. Zur Berechnung der \SourceLeaf{093}Zeitreduktion $t-T$ muss der Abstand von der Erde bekannt sein. Dafür werden meist bequeme Hilfsmittel zur Verfügung stehen, etwa eine Ephemeride, auch wenn sie nur überschlägig berechnet ist. Andernfalls genügt es, den wahren Abstand für die Zeit $t$ in einer Vorrechnung auf die übliche Weise ohne übermäßige Genauigkeit zu bestimmen.

II. Man berechne für die gegebene Zeit $t$ den wahren Ort und den Abstand, aus diesem die Zeitreduktion $t-T$ und daraus mithilfe der täglichen Bewegung (in Länge und Breite oder in Rektaszension und Deklination) die Reduktion des wahren Ortes auf die Zeit $T$.

III. Man berechne den heliozentrischen Ort der Erde für die Zeit $t$ und den heliozentrischen Ort des Planeten für die Zeit $T$. Dann bestimme man aus ihrer Verbindung auf die übliche Weise den geozentrischen Ort des Planeten. Addiert man dazu die Aberration der Fixsterne (nach dem bekannten Verfahren zu ermitteln oder aus Tafeln zu entnehmen), so erhält man den gesuchten scheinbaren Ort.

Das zweite Verfahren, das gewöhnlich angewandt wird, hat vor den übrigen den Vorteil, dass zur Bestimmung des Abstands nie eine doppelte Rechnung nötig ist. Es hat jedoch den Nachteil, dass es nur dann angewandt werden kann, wenn mehrere benachbarte Orte berechnet werden oder aus Beobachtungen schon bekannt sind. Andernfalls könnte man die tägliche Bewegung nicht als gegeben ansehen.

Der Nachteil des ersten und des dritten Verfahrens entfällt ganz, wenn mehrere benachbarte Orte zu berechnen sind. Sobald nämlich für einige die Abstände bekannt sind, lassen sich die nächstfolgenden Abstände mit den üblichen Mitteln sehr bequem und hinreichend genau ableiten. Ist der Abstand bekannt, so wird das erste Verfahren dem dritten in der Regel vorzuziehen sein, weil es die Aberration der Fixsterne nicht benötigt. Muss man hingegen auf eine doppelte Rechnung zurückgreifen, so empfiehlt sich das dritte Verfahren dadurch, dass in der zweiten Rechnung wenigstens der Ort der Erde beizubehalten ist.

Was für die umgekehrte Aufgabe erforderlich ist, ergibt sich nun von selbst, etwa wenn der wahre Ort aus dem scheinbaren hergeleitet werden soll. Nach Verfahren I behält man den Ort selbst unverändert bei, wandelt aber die Zeit $t$, der der gegebene Ort als scheinbarer entspricht, in die reduzierte Zeit $T$ um, der derselbe Ort als wahrer entspricht. Nach Verfahren II behält man die Zeit $t$ bei und addiert zum gegebenen Ort die Bewegung während der Zeit $t-T$, als wollte man ihn auf die Zeit $t+(t-T)$ reduzieren. Nach Verfahren III \SourceLeaf{094}betrachtet man den von der Aberration der Fixsterne befreiten gegebenen Ort als wahren Ort zur Zeit $T$. Der wahre Ort der Erde, der der Zeit $t$ entspricht, ist jedoch so beizubehalten, als gehörte er zur reduzierten Zeit. Der Nutzen des dritten Verfahrens wird im zweiten Buch deutlicher hervortreten.

Der Vollständigkeit halber bemerken wir noch, dass der Ort der Sonne ebenso wie der Ort eines Planeten von der Aberration betroffen ist. Da sowohl der Abstand von der Erde als auch die tägliche Bewegung nahezu konstant sind, nimmt die Aberration selbst stets einen fast konstanten Wert an, der der mittleren Bewegung der Sonne in $493^{\mathrm{s}}$ entspricht, also $=20''{,}25$. Diese Größe ist von der wahren Länge abzuziehen, um die scheinbare zu erhalten. Der genaue Wert der Aberration ist dem Produkt aus Abstand und täglicher Bewegung proportional, oder, was auf dasselbe hinausläuft, dem Abstand umgekehrt proportional. Daher wäre jener mittlere Wert im Apogäum um $0''{,}34$ zu vermindern und im Perigäum um ebenso viel zu erhöhen. Unsere Sonnentafeln enthalten bereits die konstante Aberration $-20''{,}25$. Um die wahre Länge zu erhalten, muss man daher zur tabellierten Länge $20''{,}25$ addieren.

